\section{Uniform inverse bounds for the stationary matrices}
\label{sec:ex-weighted-stationary-inverses}

We use the block Fourier norm of
Section~\ref{sec:ex-weighted-frame-penalty}. Its weight at alias $l$ is
$h_l(p)=(|q(\operatorname{Re}k_l(p))|^2+\mu^2)^{1/2}$.
The internal space is the direct sum of the fine one-form space in this
norm and the coarse scalar space in its Euclidean norm. We denote it by
$W_p$. The dual always pairs opposite momenta. The weights are not
holomorphic functions of $p$, so a complex inverse estimate requires
more than a direct application of Cauchy's inequality to a weighted
matrix.

\subsection{Continuing a real lower bound to a complex strip}
Use the maximum norm on the four momentum coordinates throughout.
Let $X_p,Y_p$ be two weighted norms, represented by invertible matrices
$S_X(p),S_Y(p)$: $\|u\|_{X_p}=|S_X(p)^{-1}u|_2$.
Suppose that, at each real momentum $x$, these matrices admit holomorphic
replacements on the ball $|z-x|<\sigma$. The replacements agree with
the actual matrices on the vertical slice $\operatorname{Re}z=x$ and
compare to the actual weighted norms in both directions by a fixed
factor $H$, and first restrict $0<\sigma\leq1$. Let $K$ be periodic and holomorphic, with weighted operator
bound $A$ on the strip of width $\sigma$, and suppose its weighted real
matrix has least singular value at least $a>0$.
Then there are $\eta>0$ and $F<\infty$, depending on $\sigma,H,A,a$
only, such that $K^{-1}:X_p\to Y_p$ has norm at most $F$ for
$p$ in the strip of width $\eta$.

Here is a quantitative proof. At a real anchor $x$, transform $K(z)$
using the holomorphic replacement matrices. The resulting matrix $Q_x(z)$
is holomorphic and has Euclidean operator norm at most $H^2A$ on the
ball of radius $\sigma$. Cauchy's estimate on complex lines, tested
against unit vectors on both sides, gives
\[
       \|Q_x(p)-Q_x(x)\|\leq k|p-x|,
             \qquad k=\frac{H^2A}{\sigma/2},
             \quad |p-x|<\sigma/2.
\]
Choose $\eta=\min\{\sigma/2,a/(2(1+k))\}$. For
$x=\operatorname{Re}p$, the actual and replacement norms coincide, and
\[
       |Q_x(p)v|\geq(a-k|p-x|)|v|\geq(a/2)|v|.
\]
Thus $Q_x(p)$ is invertible and its inverse has norm at most $2/a$.
The same is true of the weighted inverse of $K(p)$. Holomorphy of the
inverse follows from the adjugate formula; periodicity follows from
uniqueness of the inverse. No dimension enters this argument.

For the block Fourier weights the required replacement is explicit:
freeze $h_l$ at the real anchor $x$, while retaining the holomorphic
Fourier matrix $P(z)$. The inequality
\[
                    h_l(z)\leq6h_l(x),\qquad h_l(x)\leq6h_l(z)
                    \quad(|z-x|\leq1)
\]
follows from $h_l\geq\mu$, the triangle inequality and the Lipschitz
bound $|q(\operatorname{Re}k_l(z))-q(k_l(x))|\leq4\mu|z-x|$.
The Fourier matrices cancel in the comparison of the two norms; the
remaining matrix is diagonal with these ratios. Hence $H=6$ works,
uniformly in $L$. Direct sums with the ordinary scalar norm and dual
norms have the same property.

\subsection{The scalar and gauge-completed fine matrices}
Put $\omega(k)=\sum_a q_a(-k)q_a(k)$ and define
\[
 s_\mu(k)=\omega(k)+\mu^2,
 \quad d_\mu(k)=\frac{\mu^2}{s_\mu(k)},\quad
 c_\mu(k)=\mu^2d_\mu(k)^9,\quad
 \gamma_\mu(k)=\omega(k)+c_\mu(k).
\]
At real $k$, $\omega(k)=|q(k)|^2$. The sign convention in this display
uses $q_a(-k)q_a(k)=|q_a(k)|^2$ on the real slice.
For $x\geq0$,
\begin{equation}\label{eq:ex-scalar-weight-floor}
 x+\mu^2\left(\frac{\mu^2}{x+\mu^2}\right)^9
                   \geq2^{-10}(x+\mu^2).
\end{equation}
If $x\geq\mu^2$, its first term alone suffices. If $x<\mu^2$,
the ratio is greater than $1/2$, so the second term is at least
$2^{-9}\mu^2\geq2^{-10}(x+\mu^2)$.
On the alias strip $|\operatorname{Im}k|_\infty<\mu/128$,
comparison with the real symbol gives
$|s_\mu(k)|\geq\tfrac12(|q(\operatorname{Re}k)|^2+\mu^2)$.
Indeed the change of $q$ has norm at most $4\mu/128$ and the change
of its quadratic product is bounded by twice this times the real norm
plus its square; $2ab\leq a^2+b^2$ absorbs this into less than half
the real denominator. In particular all displayed rational symbols are
holomorphic there, $|d_\mu|\leq2$, and
$|c_\mu|\leq2^9\mu^2$.

Fix $\varepsilon>0$, $\Theta\geq0$ and a coarse width
$0<\sigma\leq1$. Let $H(k)$ be the symbol of the fine two-form metric,
holomorphic and bounded by $\Theta$ on the fine strip
$|\operatorname{Im}k|_\infty<\mu\sigma$, which contains the aliases
$k_l(p)$ of the coarse strip of width $\sigma$. Write
$K_f=d_1(-k)^TH(k)d_1(k)$ and $d_0(k)v=q(k)\wedge v$, and suppose that
the real kinetic form satisfies the all-period lower bound
\[
        \langle v,K_f(k)v\rangle\geq\varepsilon|d_1(k)v|^2
\]
at every momentum $k$ of every finite-period Fourier grid, with the same
$\varepsilon$ for all periods. In the application this is a consequence
of the floor condition of Section~\ref{sec:ex-soft-curvature-receipts},
which gives $\langle w,H_mw\rangle\geq\varepsilon|w|^2$ for every real
two-form $w$ at every period $m$. The positivity of $\varepsilon$ is
used twice: the weighted floor of the gauge-completed matrix $D_\mu$
and the constant $a_0$ of the quotient bound
\eqref{eq:ex-quotient-floor} below are positive only because
$\varepsilon>0$.
The scalar matrix $s_\mu I$ and the scalar matrix $\gamma_\mu I$
have weighted real floors $1$ and $2^{-10}$. The gauge-completed
one-form matrix
\[
       D_\mu(k)=K_f(k)+d_0(k)d_0(-k)^T+c_\mu(k)I
\]
has weighted floor $2^{-10}\min\{\varepsilon,1\}$.
This follows from
\[
       |d_1v|^2+|d_0^*v|^2=|q|^2|v|^2
\]
and \eqref{eq:ex-scalar-weight-floor}. The exterior identity is obtained
by expanding the wedge and contraction in the coordinate basis.
The all-period kinetic bound applies at every real momentum by density
of finite-torus frequencies and continuity.

For these scalar and fine symbols restrict further to $\sigma\leq1/128$.
Their complex weighted upper bounds are also uniform. Each exterior
matrix costs at most a fixed color-direction constant times $|q|$,
$|q(k_l)|\leq5h_l$, and $|c_\mu(k_l)|\leq2^9h_l^2$.
Thus division by $h_l$ at both ends leaves a fixed bound depending only
on $G,\Theta$. The block Fourier conjugations cancel between opposite
momenta, leaving a unitary alias permutation and these diagonal weights.
The preceding continuation argument gives inverse bounds for all three
matrices on one coarse strip, before $L$ and before any higher-order
interaction budget are chosen.

\subsection{A real lower bound modulo gauge directions}
Let $Q(p)$ be the linear block mean on one-forms and $J(p)$ its scalar
counterpart, so $Qd_0=d_{0,c}J$. Set
\[
 \mathcal R a=(-d_0a,-Ja),\qquad
 N(u,r)=d_{0,c}r-Qu,
 \qquad \mathcal G(u,r)=\langle u,K_fu\rangle+\kappa|N(u,r)|^2,
\]
where $\kappa>0$ and all operators are expressed in the same block
Fourier coordinates. The curl and $N$ annihilate $\mathcal R a$, so
$\mathcal G$ is unchanged when a gauge vector is subtracted.
Recall that $\varepsilon>0$ and $\kappa>0$. We claim that, for some
$a_0>0$ depending only on $G,J,\kappa,\varepsilon$,
\begin{equation}\label{eq:ex-quotient-floor}
        \text{for every }v\text{ there is }a\text{ such that }
        a_0\|v-\mathcal R a\|_{W_p}^2\leq\mathcal G(v).
\end{equation}

Choose the real coarse momentum in $[-\pi,\pi]^4$. For every nonzero
alias, $|q(k_l)|\geq2\mu$. Subtract its longitudinal component by
solving $d_0(k_l)^*(u_l+d_0(k_l)a_l)=0$; the scalar denominator is
$|q(k_l)|^2>0$. The principal alias is left available to eliminate the
coarse scalar coordinate. In normalized Fourier coordinates the scalar
block mean on that alias is $\mu^2\operatorname{rollMean}(k_0)$,
which is nonzero. No uniform bound on the eliminating gauge parameter
is needed. Here $\operatorname{rollMean}$ is the product
of the four normalized half-block geometric sums, raised to power $J$.
For $|k_a|\leq\pi/L$, each sum has modulus at least $1/16$;
this follows from the geometric-sum formula and
$\sin t\geq2t/\pi$ for $0\leq t\leq\pi/2$, with the value at zero
understood by continuity. Thus $|\operatorname{rollMean}|\geq c_J=(1/16)^{4J}$. We obtain a representative $(P(p)b,0)$ with all
nonprincipal $b_l$ transverse.

Put $A_h=\sum_{l\ne0}h_l^2|b_l|^2$ and
$B_0=|QP(p)b|^2$. The kinetic bound and transversality give
\[
             \mathcal G(v)\geq(\varepsilon/2)A_h+\kappa B_0.
\]
The principal part of $QP(p)$ has smallest singular value at least
$(c_J/16)\mu$: the straight connector contributes a geometric sum of
size at least $L/16$, and $P$ contributes $\mu^2$.
The full block mean has weighted norm at most a constant $C_J$.
To verify this upper bound directly, each fixed-position read has
coefficients bounded by $C_J\mu^4$ on $O_J(L^4)$ sites, hence Euclidean
row norm at most $C'_J\mu^2$; sum the $L$ positions and use
$h_l\geq\mu$. The same argument gives the scalar averaging norm
$\|J\|\leq C'_J\mu^2$.
Separating the principal and nonprincipal aliases and squaring the
triangle inequality now gives
\[
 ((c_J/16)\mu|b_0|)^2\leq2B_0+2C_J^2A_h.
\]
Finally $h_0^2\leq((2\pi)^2+1)\mu^2$. It follows that
$\|(P b,0)\|_{W_p}^2=A_h+h_0^2|b_0|^2$ is bounded by a fixed
multiple of $(\varepsilon/2)A_h+\kappa B_0$, proving
\eqref{eq:ex-quotient-floor}. This proof includes $p=0$.

\subsection{Adding the gauge-fixing square}
Let $F$ be the actual linear gauge-fixing map, periodic and holomorphic
on the chosen strip, and satisfying $F(-p)^T=F(p)^*$ on real momenta.
Its required identities
and estimate are
\begin{equation}\label{eq:ex-gauge-fixing-receipt}
       F\mathcal R=\Gamma_\mu,
       \qquad |Fv|\leq D\|v\|_{W_p},
\end{equation}
where the bound holds throughout that complex strip, and $\Gamma_\mu$
is the block scalar operator with symbol $\gamma_\mu$. These are exact same-frame identities; they remain
inputs to this subsection until the gauge-fixing construction is
identified with these symbols.
The bound on $J$, the definition of $\mathcal R$ and
\eqref{eq:ex-scalar-weight-floor} give
\[
       \|\mathcal R a\|_{W_p}\leq C\,|\Gamma_\mu a|
                                 =C|F\mathcal R a|,
\]
with $C$ independent of $L$. For the fine component this is the
pointwise inequality $h_l|q_l|\leq h_l^2\leq2^{10}\gamma_\mu(k_l)$;
for the coarse component use $\mu^2\leq h_l^2$ and the scalar mean
bound above.
Given $v$, take $a$ from \eqref{eq:ex-quotient-floor} and put
$b=\|v-\mathcal R a\|_{W_p}$, $c=|Fv|$. Then
\[
       \|v\|_{W_p}\leq Cc+(CD+1)b,
       \qquad a_0b^2+c^2\leq\mathcal G(v)+|Fv|^2.
\]
Squaring the first inequality proves a uniform positive lower bound
for the full quadratic form $\mathcal G+F^*F$; for example its inverse
floor can be taken as
$\max\{2C^2,2(CD+1)^2/a_0\}$.
On the complex strip its weighted upper bound is the sum of the
kinetic bound, $\kappa$ times the square of the bound for $N$, and $D^2$.
The holomorphic opposite-momentum transpose replaces $F^*$ off the real
slice. The continuation lemma therefore gives a common inverse width
and cost for this actual stationary fluctuation matrix.

Once the common matrix is identified on the Fourier grids of an
unbounded family of tori, the identification extends to the entire
real slice by density, and then to the strip by analytic uniqueness.
Consequently any two stationary matrix presentations representing that
same family inherit this one inverse estimate. This last statement does
not identify arbitrary independently chosen matrices: the all-period
representation and the exact identities
\eqref{eq:ex-gauge-fixing-receipt} must refer to the same construction.
